% Môn: Vật lý
% Lớp: 11
% Chương: Chương I. Dao động
% Bài: L11C1 Bài 4. Bài tập về dao động điều hòa · Viết phương trình dao động điều hòa từ đồ thị li độ – thời gian
% ===== Câu 20 =====
% ID: L11C1B4-33-DS
% Mức: TH
\dangbt{Viết phương trình dao động điều hòa từ đồ thị li độ – thời gian}
\begin{ex}
% ID: L11C1B4-33-DS
% Mức: TH
Khi phân tích các đặc điểm của đồ thị li độ – thời gian của một vật dao động điều hòa quanh vị trí cân bằng, hãy nhận xét các mệnh đề sau:
\choiceTF
{\True Tọa độ li độ $x$ của vật tại một thời điểm được xác định trực tiếp bằng tung độ của điểm biểu diễn tương ứng trên đồ thị.}
{\True Độ dốc của đồ thị tại một điểm bất kì biểu thị giá trị vận tốc tức thời của vật tại thời điểm đó.}
{\True Nếu đường hình sin của đồ thị có xu hướng đi xuống cắt trục hoành thì vật đang chuyển động qua vị trí cân bằng theo chiều âm.}
{Khoảng cách giữa hai điểm cắt trục hoành liên tiếp của đồ thị luôn có giá trị bằng một chu kì dao động $T$.}
\loigiai{
\begin{itemchoice}
\item (Đúng) Tọa độ li độ $x$ của vật tại một thời điểm được xác định trực tiếp bằng tung độ của điểm biểu diễn tương ứng trên đồ thị. (Vì): Mệnh đề thứ nhất đúng vì trục tung biểu diễn li độ $x$, do đó giá trị tung độ chính là tọa độ li độ thực tế của vật
\item (Đúng) Độ dốc của đồ thị tại một điểm bất kì biểu thị giá trị vận tốc tức thời của vật tại thời điểm đó. (Vì): Mệnh đề thứ hai đúng vì theo định nghĩa vật lí, vận tốc là đạo hàm của li độ theo thời gian, tương ứng với độ dốc của tiếp tuyến đồ thị $x - t$
\item (Đúng) Nếu đường hình sin của đồ thị có xu hướng đi xuống cắt trục hoành thì vật đang chuyển động qua vị trí cân bằng theo chiều âm. (Vì): Mệnh đề thứ ba đúng vì khi đồ thị cắt trục hoành ($x = 0$) đi xuống nghĩa là li độ đang giảm dần về giá trị âm, tức là vật chuyển động ngược chiều dương
\item (Sai) Khoảng cách giữa hai điểm cắt trục hoành liên tiếp của đồ thị luôn có giá trị bằng một chu kì dao động $T$. (Vì): Mệnh đề thứ tư sai vì hai điểm cắt trục hoành liên tiếp tương ứng với hai lần liên tiếp vật đi qua vị trí cân bằng, khoảng thời gian này chỉ bằng một nửa chu kì $\dfrac{T}{2}$
\end{itemchoice}
}
\end{ex}

% ===== Câu 21 =====
% ID: L11C1B4-34-DS
% Mức: VD
\begin{ex}
% ID: L11C1B4-34-DS
% Mức: VD
Quan sát đồ thị dao động điều hòa của một vật có đường hình sin biểu diễn trên hệ tọa độ $x - t$. Tại thời điểm $t = 0$, đồ thị đi qua gốc tọa độ và đang đi lên theo chiều dương, sau đó đạt tới biên dương với tọa độ li độ $10$ cm lần đầu tiên tại mốc thời gian $t = 0,25$ s. Khảo sát các mệnh đề sau:
\choiceTF
{\True Biên độ dao động của vật là $A = 10$ cm.}
{\True Chu kì dao động của vật là $T = 1,0$ s.}
{Pha ban đầu của dao động bằng $\varphi = \dfrac{\pi}{2}$ rad.}
{\True Phương trình li độ của vật là $x = 10\cos\left(2\pi t - \dfrac{\pi}{2}\right)$ (cm).}
\loigiai{
\begin{itemchoice}
\item (Đúng) Biên độ dao động của vật là $A = 10$ cm. (Vì): Mệnh đề thứ nhất đúng vì tọa độ đỉnh của đường hình sin trên đồ thị tương ứng với biên độ $A = 10$ cm
\item (Đúng) Chu kì dao động của vật là $T = 1,0$ s. (Vì): Mệnh đề thứ hai đúng vì khoảng thời gian vật đi từ vị trí cân bằng đến vị trí biên dương lần đầu tiên bằng một phần tư chu kì: $\dfrac{T}{4} = 0,25 \text{ s} \Rightarrow T = 1,0$ s
\item (Sai) Pha ban đầu của dao động bằng $\varphi = \dfrac{\pi}{2}$ rad. (Vì): Mệnh đề thứ ba sai vì tại gốc thời gian vật đi qua vị trí cân bằng theo chiều dương nên pha ban đầu có giá trị âm: $\varphi = -\dfrac{\pi}{2}$ rad
\item (Đúng) Phương trình li độ của vật là $x = 10\cos\left(2\pi t - \dfrac{\pi}{2}\right)$ (cm). (Vì): Mệnh đề thứ tư đúng vì tần số góc $\omega = \dfrac{2\pi}{T} = 2\pi$ rad/s, kết hợp với biên độ và pha ban đầu ta viết được phương trình chuẩn
\end{itemchoice}
}
\end{ex}

% ===== Câu 23 =====
% ID: L11C1B4-35-DS
% Mức: VD
\begin{ex}
% ID: L11C1B4-35-DS
% Mức: VD
Một vật dao động điều hòa có đồ thị li độ – thời gian dạng hình sin. Tại thời điểm $t = 0$, đồ thị bắt đầu từ biên dương $x_0 = 5$ cm và đi xuống vị trí cân bằng lần đầu tiên tại thời điểm $t = 0,1$ s. Xét các thông số của dao động:
\choiceTF
{\True Biên độ dao động của vật bằng $5$ cm.}
{\True Chu kì dao động của vật bằng $0,4$ s.}
{\True Tần số góc của dao động đạt giá trị $5\pi$ rad/s.}
{Phương trình dao động điều hòa của vật là $x = 5\cos\left(5\pi t + \dfrac{\pi}{2}\right)$ (cm).}
\loigiai{
\begin{itemchoice}
\item (Đúng) Biên độ dao động của vật bằng $5$ cm. (Vì): Mệnh đề thứ nhất đúng vì giá trị li độ lớn nhất ở thời điểm ban đầu là biên độ $A = 5$ cm
\item (Đúng) Chu kì dao động của vật bằng $0,4$ s. (Vì): Mệnh đề thứ hai đúng vì thời gian vật đi từ vị trí biên về vị trí cân bằng lần đầu tiên là một phần tư chu kì: $\dfrac{T}{4} = 0,1 \text{ s} \Rightarrow T = 0,4$ s
\item (Đúng) Tần số góc của dao động đạt giá trị $5\pi$ rad/s. (Vì): Mệnh đề thứ ba đúng vì tần số góc tương ứng là: $\omega = \dfrac{2\pi}{T} = \dfrac{2\pi}{0,4} = 5\pi$ rad/s
\item (Sai) Phương trình dao động điều hòa của vật là $x = 5\cos\left(5\pi t + \dfrac{\pi}{2}\right)$ (cm). (Vì): Mệnh đề thứ tư sai vì tại $t = 0$ vật ở biên dương nên pha ban đầu $\varphi = 0$. Phương trình đúng phải là: $x = 5\cos(5\pi t)$
\end{itemchoice}
}
\end{ex}

% ===== Câu 24 =====
% ID: L11C1B4-36-TLN
% Mức: VD
\begin{ex}
% ID: L11C1B4-36-TLN
% Mức: VD
Trên đồ thị li độ–thời gian, hai đỉnh dương liên tiếp ở $t=1$ s và $t=5$ s. Chu kì theo đơn vị s bằng bao nhiêu?
\shortans{03/01/1900}
\end{ex}

% ===== Câu 37 =====
% ID: L11C1B4-37-TN
% Mức: NB
\begin{ex}
% ID: L11C1B4-37-TN
% Mức: NB
$\immini$ {Dựa vào đồ thị li độ – thời gian như hình vẽ, biên độ dao động của vật là:
\choice
{$10\ \mathrm{cm}$}
{$15\ \mathrm{cm}$}
{\True $20\ \mathrm{cm}$}
{$25\ \mathrm{cm}$}
\loigiai{
Biên độ là độ lớn cực đại của li độ , $A = 20\ \mathrm{cm}$.
}
\end{ex}

% ===== Câu 39 =====
% ID: L11C1B4-38-TN
% Mức: VD
\begin{ex}
% ID: L11C1B4-38-TN
% Mức: VD
Đồ thị li độ – thời gian cho thấy vật có biên độ $4\,\text{cm}$, chu kì $2\,\text{s}$; tại t = 0 vật qua vị trí cân bằng theo chiều dương. Phương trình dao động nào phù hợp?
\choice
{x = 4cos(πt + π/2) cm}
{\True x = 4cos(πt - π/2) cm}
{x = 4cos(2πt) cm}
{x = 4cos(2πt - π/2) cm}
\loigiai{
$T=2\,\text{s}$ nên ω = π rad/s. Tại t = 0, x = 0 và v > 0 nên pha ban đầu là -π/2.
}
\end{ex}

% ===== Câu 42 =====
% ID: L11C1B4-39-TN
% Mức: VD
\begin{ex}
% ID: L11C1B4-39-TN
% Mức: VD
Đồ thị li độ – thời gian cho thấy vật có biên độ $6\,\text{cm}$, chu kì $4\,\text{s}$ và ở biên dương tại t = 0. Phương trình dao động nào phù hợp?
\choice
{\True x = 6cos(πt/2) cm}
{x = 6cos(2πt) cm}
{x = 6cos(πt/2 + π/2) cm}
{x = 6cos(πt) cm}
\loigiai{
$T=4\,\text{s}$ nên ω = π/2 rad/s. Vật ở biên dương lúc t = 0 nên pha ban đầu bằng 0.
}
\end{ex}
